\chapter{Conventions}
\label{app:conventions}

\section{Grading conventions}
\label{sec:grading-conventions}

All grading is \textbf{cohomological}: differentials have degree $+1$.  The bar construction uses desuspension (matching Volumes~I--II).

\section{Sign conventions}
\label{sec:sign-conventions}

Signs follow the Koszul rule throughout.  For the curved $\Ainf$-structure: $\mu_1^2(a) = [\mu_0, a]$ (commutator, minus sign).

\section{Categorical conventions}
\label{sec:categorical-conventions}

\begin{itemize}
  \item $\DGCat$ denotes the $\infty$-category of small dg categories over $k$.
  \item $\CY_d\text{-}\Cat$ denotes the $\infty$-category of smooth $d$-dimensional CY categories.
  \item Tensor products of dg categories are derived unless stated otherwise.
  \item ``Equivalence'' means quasi-equivalence of dg categories (equivalence of associated $\infty$-categories).
\end{itemize}

\section{Operadic conventions}
\label{sec:operadic-conventions}

\begin{itemize}
  \item $E_n$ denotes the little $n$-disks operad (Boardman--Vogt).
  \item $\FM_k(\mathbb{C})$ denotes the Fulton--MacPherson compactification of $\mathrm{Conf}_k(\mathbb{C})$.
  \item $\SC^{\mathrm{ch,top}}$ denotes the Swiss-cheese operad from Volume~II.
\end{itemize}
